\bookchapter{Chapter 33: Future Outlook — The Next Journey of Smartphones}

\chapterintro{Phones won't disappear overnight, but their “monopoly” may be eroded by a new generation of devices. This chapter covers three things: what next-gen devices look like, how the phone may “disappear” to become a new identity foundation, and the ethical issues we must consider ahead of time.}

\booksection{I. AR/VR and Spatial Computing: The “Second Desktop” Beyond Phones}

The competition for next-gen devices has begun. The most likely candidates to divert time from phones are AR/VR and spatial computing devices:

\begin{itemize}
  \item \textbf{Vision Pro Successor}: The launch of Apple Vision Pro marked the official arrival of the "spatial computing" concept, but the first generation's price (starting at 29,999 yuan) and weight remain barriers. A lighter, cheaper Vision is expected in 2026-2027, bringing spatial computing from "developer toy" to "early majority." At that point, iPhone may become the computing companion and data hub for Vision glasses.
  \item \textbf{Another future evolution direction for phones may be "invisibility." It will transform into an omnipresent "behind-the-scenes promoter," becoming the sole cornerstone of your digital identity.}: Open a "second desktop" in front of your eyes, making digital information as ubiquitous as air. Meta and Ray-Ban's smart glasses have sold millions of pairs, and domestic brands like Rokid, XREAL, and INMO are also releasing consumer AR glasses. By 2025-2026, AR glasses can do real-time translation, navigation HUD, and AI assistant conversations, but they're still limited by waveguide display and battery life — 3-5 years from "replacing phones."
  \item \textbf{Smart Rings and Wearables}: Samsung Galaxy Ring, Oura Ring, and other smart rings are gaining popularity, serving as "body-worn sensors" that complement the phone's health monitoring. The phone's role is shifting from "all-in-one terminal" to "hub + multiple lightweight accessories."
\end{itemize}

\booksection{II. Brain-Computer Interface: Thought as Interaction Medium}

Thought itself will become an interaction medium, but this path is longer than imagined:

\begin{itemize}
  \item \textbf{Neuralink Latest Progress}: Musk's Neuralink completed its first human implant in 2024 — the first subject, Noland Arbaugh, controlled a computer mouse with his thoughts to play chess. Currently, implantable BCIs (brain-computer interfaces) primarily target assistive scenarios for paralyzed patients, aiming to help them "operate" the digital world again.
  \item \textbf{Non-invasive BCI}: Multiple companies are developing non-invasive headbands (such as Kernel, OpenBCI) that can read brain signals without surgery, but with far lower precision than invasive ones. We're still quite far from "typing with thoughts."
  \item \textbf{Consumer Outlook}: BCI is 10-20 years from consumer-level phone interaction, but may first be used in accessibility assistance and focus monitoring. Ethical and privacy issues ("Whose privacy do my brain waves belong to?") will become debate focal points before the technology matures.
\end{itemize}

\booksection{III. The Phone's “Disappearance” and New Identity Foundation}

Another future evolution direction for the phone may be “disappearing” — becoming an omnipresent “behind-the-scenes enabler” and the sole foundation of your digital identity:

\begin{itemize}
  \item \textbf{: Your cloud personal data vault is the center.}: Instantly complete "prove you are you" authentication — no need to manually enter verification codes when logging into any service.
  \item \textbf{An ultimate question always exists: Are we "gods" empowered by technology, or "brains in vats" fed by algorithms?}: Pre-set digital asset inheritance plans during new device activation, making "digital afterlife" a factory default.
  \item \textbf{When you unconsciously light up the screen every few minutes, is this your habit or a designed "addictive behavior"?}: Your cloud-based personal data vault is the center; hardware is just an access point. Change phones or devices — your identity and data always follow you.
\end{itemize}

\booksection{IV. Ethical Reflection: Are We Still the Masters of Our Phones?}

Let automation be your faithful “co-pilot.”

\begin{itemize}
  \item \textbf{AI Ethics}: When on-device AI understands your preferences better than you do, is your "choice" an expression of free will or a guided outcome? When AI Agents place orders and reply to messages on your behalf, where do responsibility and boundaries lie?
  \item \textbf{Data Sovereignty}: Does your health data, brain wave data, and biometric data belong to you, the device manufacturer, or the cloud service provider? The Personal Information Protection Law and EU GDPR are just starting points — the true implementation of "data portability rights" and "the right to be forgotten" still has a long way to go.
  \item \textbf{Digital Divide}: AI phones, foldable screens, and satellite communications are widening the gap between those who "can afford" and those who "can't." Equitable access to digital infrastructure will become a new social issue — just as critical as "knowing how to use the internet" was 20 years ago.
  \item Build your stable "digital rear base."
  \item Regularly perform "digital decluttering" and "social fasting."
\end{itemize}

\booksection{V. Stay Connected, Stay Independent — Your Digital Life Philosophy}

No matter how technology evolves, these digital life philosophies will still apply 10 years from now:

\begin{enumerate}
  \item Try It: Draft Your "Personal Digital Device Terms of Service"
  \item Write a simple "terms of service" for you and your phone to define your relationship. For example:
  \item Article 1: My attention time belongs to me. You should proactively enter "Do Not Disturb" mode when I am working/spending time with family.
  \item Article 2: My personal data is my property. Your primary responsibility is secure encryption and convenient syncing.
  \item Article 3: Our common goal is to improve life efficiency and experience. If you make me addicted through manipulative design, I have the right to uninstall you.
\end{enumerate}

\begin{tcbpractice}{V. Appendices}
Write a simple “terms of service” for yourself and your phone (and future AR glasses and AI assistants) to define your relationship. For example:

\begin{itemize}
  \item Appendix A: New Phone Unboxing "Step-by-Step" Guide
  \item 1. Unboxing Inspection "Three-Step Process":
  \item Check Appearance
\end{itemize}

\end{tcbpractice}

\begin{tcbsummary}{Chapter Summary}
\begin{itemize}
  \item Vision Pro and AR glasses are new computing entry points beyond phones, but won't replace phones in the short term
  \item BCI is 10-20 years from consumer use; ethical issues will mature before the technology does
  \item The phone is evolving toward "disappearing" — from device-centric to identity-centric
  \item AI ethics, data sovereignty, and the digital divide are core issues for the next 10 years
  \item Staying connected while staying independent is the timeless philosophy of digital life
\end{itemize}

\end{tcbsummary}
